% Text of the poster. The figures come from figures/, made by poster_figures.py.
\newcommand{\TitleLines}{Meshing smooth CSG\\with Lipschitz pruning}
\newcommand{\Subtitle}{Adaptive dual contouring where every octree cell evaluates only its
  pruned tree}
\newcommand{\Course}{Computer Graphics}
\newcommand{\Date}{November 2026}
\newcommand{\TeaserA}{Trees, 371 nodes}
\newcommand{\TeaserB}{Monument, 6323 nodes}
\newcommand{\TeaserC}{Molecule, 2001 nodes}
\newcommand{\ScaleText}{Active nodes of the pruned tree at each vertex, from 1 to 20}

\newcommand{\SecIntro}{Introduction}
\newcommand{\SecContribution}{Contribution}
\newcommand{\SecResults}{Results}
\newcommand{\SecLimitations}{Limitations}
\newcommand{\SecConclusions}{Conclusions}
\newcommand{\SecReferences}{References}

% ---------------------------------------------------------------- introduction
\newcommand{\IntroText}{%
  A signed distance function (SDF) can be built as a tree whose leaves are primitives and whose
  internal nodes are hard or smooth CSG operators with a blending radius \(k\) [5]. To render it,
  sphere tracing [4] marches along each ray and evaluates the SDF at every step. Each evaluation
  traverses the whole tree, so the cost grows with the number of primitives.\par
  \vspace{2mm}
  Barbier et al.\ [1] compute local pruned trees that are equivalent to the full tree within
  their region of space and much faster to evaluate. For each cell of a grid they keep the
  active nodes, those needed for at least one point of the cell. Thanks to the Lipschitz
  property of SDFs, one evaluation of the tree at the center of the cell is enough to tell which nodes are
  unnecessary in the whole cell. The pruning is hierarchical, since the grid is refined from
  coarse to fine and each cell prunes the already pruned tree of its parent cell.\par
  \vspace{2mm}
  In the cells far from the surface, where \(|f(c)|\) is large compared with \(R\) (\(c\) being the
  center of the cell and \(R\) the radius of its sphere), they store no tree, only the constant
  \(|f(c)| - R\) with the sign of \(f(c)\). By the Lipschitz property, that constant is still
  a lower bound on the distance in the whole cell. They call this far-field culling. They apply the method to sphere tracing and to discretizing the SDF on a dense grid, and note that it
  could also be used for polygonization.\par}

% ---------------------------------------------------------------- contribution
\newcommand{\ContributionText}{%
  The sparse octree implemented here subdivides only the cells that may
  contain the surface and gives each one a pruned tree computed from that of its parent cell.
  In addition, for dynamic scenes, moving a primitive re-prunes only the cells it reaches. The paper leaves
  this sparse partition and the partial update as future work, and both were also
  implemented on the GPU with Vulkan.\par
  \vspace{2mm}
  The paper only suggests polygonization. On the leaves of the octree, dual contouring [2] builds
  the mesh, and in dynamic scenes only the contour of the cells that
  changed is rebuilt. Where the surface is simple, as on
  a flat face, the mesh uses larger cells. While the octree is built, a cell whose pruned tree
  has a single primitive stops being subdivided if one vertex is enough to approximate the
  surface inside it. Dual contouring and the contour patch
  run on the GPU as well, and a dedicated viewer draws the mesh without leaving the GPU.\par}

% ---------------------------------------------------------------- results
\newcommand{\NumA}{228\(\times\)}
\newcommand{\NumAText}{faster than the full tree when meshing Monument on the GPU.}
\newcommand{\NumB}{0}
\newcommand{\NumBText}{differences between the meshes of the full tree and Lipschitz pruning.}
\newcommand{\NumC}{123}
\newcommand{\NumCText}{frames per second dragging a sphere through Monument.}

\newcommand{\FigTimes}[1]{%
  \includegraphics[width=#1mm]{times.pdf}\par
  {\captiontext Warm GPU time at depth 8, on a logarithmic scale, with the mesh kept on the GPU. The
   {\bfseries full tree}, \textcolor{interval}{\bfseries intervals} and \textcolor{method}{\bfseries Lipschitz}
   produce the same mesh, and the parentheses give the speedup of Lipschitz over the full tree. The dynamic
   frame is the partial update when a sphere moves.\par}}

\newcommand{\FigDynamic}[1]{%
  \includegraphics[width=#1mm]{dynamic.png}\par
  \vspace{2mm}
  {\captiontext One frame of the sphere crossing Monument. The patch rebuilds only the triangles in
   \textcolor{method}{\bfseries burgundy} and keeps the rest of the mesh.\par}}

\newcommand{\FigCoarseDetail}{%
  \begin{minipage}[c]{58mm}\centering\includegraphics[width=58mm]{coarse_overview.png}\par\vspace{1.5mm}%
    {\captiontext part 145\par}\end{minipage}\hfill
  \begin{minipage}[c]{96mm}\raggedright
    \includegraphics[width=96mm]{coarse_uniform_zoom.png}\par\vspace{1mm}{\captiontext uniform mesh, 791 k triangles\par}
    \vspace{3mm}
    \includegraphics[width=96mm]{coarse_coarse_zoom.png}\par\vspace{1mm}{\captiontext coarser leaves, 304 k triangles\par}
  \end{minipage}\par}

\newcommand{\ResultsTable}{%
  {\tablefont
   \begin{tabular}{@{}lrrr@{}}
     \tabhead{depth 8, GPU} & \tabhead{Trees} & \tabhead{Monument} & \tabhead{Molecule}\\[1mm]
     triangles & 429 k & 1.18 M & 504 k\\
     active nodes per leaf & 5.8 & 5.6 & 7.5\\
     leaves with one primitive & 51~\% & 37~\% & 13~\%\\
     pruned trees, octree & 36 MB & 149 MB & 93 MB\\
     pruned trees, grid of the paper & 83 MB & 294 MB & 198 MB\\
   \end{tabular}\par}}

\newcommand{\FigScaling}[1]{%
  \includegraphics[width=#1mm]{scaling.pdf}\par
  {\captiontext Random spheres of mean radius \(r = 0.05\) in \([-1, 1]^3\), with \(k = 0\) solid and
   \(k = r\) dashed, for the {\bfseries full tree},
   \textcolor{interval}{\bfseries intervals} and \textcolor{method}{\bfseries Lipschitz}.\par}}

\newcommand{\FigCoarse}[1]{%
  {\tablefont
   \begin{tabular*}{#1mm}{@{\extracolsep{\fill}}lrrrrrr@{}}
     \tabhead{coarser leaves} & \tabhead{lock} & \tabhead{159} & \tabhead{161} & \tabhead{145} & \tabhead{base} & \tabhead{mount}\\[1mm]
     triangles, uniform & 867 k & 1.00 M & 543 k & 791 k & 474 k & 418 k\\
     triangles, coarse & 217 k & 326 k & 179 k & 304 k & 183 k & 190 k\\
     time, uniform & 14 ms & 11 ms & 10 ms & 19 ms & 11 ms & 10 ms\\
     time, coarse & 18 ms & 29 ms & 17 ms & 32 ms & 23 ms & 16 ms\\
   \end{tabular*}\par}
  \vspace{3mm}
  {\captiontext CAD parts from InverseCSG [6] at depth 8, on the GPU, with the detail of part 145 above. The meshes
   are closed, and their largest distance to the surface is that of the uniform mesh, except on 161
   (0.19~cells versus 0.05).\par}}

% ---------------------------------------------------------------- limitations
\newcommand{\LimitationsText}{%
  The original method only supports 1-Lipschitz primitives and operators, and prunes less as the
  blending radius grows, until the cost approaches that of the full tree. Densely packed
  primitives have the same effect.\par
  \vspace{2mm}
  In the new method, dual contouring leaves non-manifold edges, 3 on Molecule and 19 on
  Monument. Coarser leaves barely pay off under smooth blending, and remove no triangles on
  Molecule. On the GPU they also take up to two and a half times as long as the uniform mesh. The GPU version
  reaches depth 9 at most, because of the dense grids of dual contouring. Byte-identical
  meshes force it to use double precision, which on the GPU used is 64 times slower than single
  precision.\par}

% ---------------------------------------------------------------- conclusions
\newcommand{\ConclusionsText}{%
  The pruning carries over to meshing exactly, and the sparse octree replaces the dense grid of
  the paper. On Monument the pruned trees of the octree take 149~MB, versus 294~MB on the
  \(257^3\) grid of the paper. Compared with interval arithmetic [3], Lipschitz pruning meshes
  Trees, Monument and Molecule 1.8, 1.5 and 2.2 times faster. On hard
  CSG, intervals evaluate somewhat fewer, but where they run out of memory, Lipschitz degrades slowly.\par
  \vspace{2mm}
  On this basis, coarser leaves remove 40 to 80~\% of the triangles on CAD parts [6], and on the
  GPU, with a sphere crossing Monument, the scene is re-meshed from scratch in 107~ms per frame. Keeping the 97~\% of
  cells it does not reach and patching only the contour brings the frame down to 6.8~ms, with a
  mesh byte-identical to the one rebuilt from scratch.\par}

% ---------------------------------------------------------------- references
\newcommand{\ReferencesText}{%
  \bib{1}W.~Barbier, M.~Sanchez, A.~Paris, É.~Michel, T.~Lambert, T.~Boubekeur, M.~Paulin,
  T.~Thonat. Lipschitz pruning: hierarchical simplification of primitive-based SDFs.
  \emph{Computer Graphics Forum} 44(2), 2025.\par
  \bib{2}T.~Ju, F.~Losasso, S.~Schaefer, J.~Warren. Dual contouring of Hermite data.
  \emph{ACM TOG} 21(3), 2002.\par
  \bib{3}M.~J.~Keeter. Massively parallel rendering of complex closed-form implicit surfaces.
  \emph{ACM TOG} 39(4), 2020. libfive.com\par
  \bib{4}J.~C.~Hart. Sphere tracing: a geometric method for the antialiased ray tracing of
  implicit surfaces. \emph{The Visual Computer} 12(10), 1996.\par
  \bib{5}I.~Quílez. Smooth minimum: normalization, thickness and bounds. iquilezles.org,
  2024.\par
  \bib{6}T.~Du et al. InverseCSG: automatic conversion of 3D models to CSG trees. \emph{ACM
  TOG} 37(6), 2018.\par}
\newcommand{\Credits}{The scenes Trees and Monument are by É.~Michel (CC-BY 4.0) and Molecule is from Galin
  et al.\ 2020. Measurements by the author.}
